Alan, this is where I stand on this ... August 20, 2000
Jim Wednesday 6:30pm

Edits for Case Study 1: Eddy Current


-----Section 0: Problem Description-----

Background & Project Goal:
The goal of this project is to develop a
nondestructive portable device for detecting
cracks and fractures in metals.  A primary
application for such would be for the detection of
defects in airplane wings.  The internal mechanism
for the detector would be to sense crack-induced
changes in the detector's electromagnetic field
which would in turn result in changes in the
impedance level of the detector.  This change of
impedance is termed "sensitivity" and it is a
sub-goal of this experiment to maximize such
sensitivity as the detector is moved from an
unflawed region to a flawed region on the metal.

Statistical Goals:
The specific statistical goals of the experiment
are the usual:

   1. Factors : determine important factors which
                affect sensitivity;
   2. Settings: determine settings which
                maximize sensitivity;
   3. Model   : determine a prediction equation which
                functionally relates sensitivity
                to the various factors.

There were 3 detector wiring component factors
under consideration:

   X1 = Number of wire turns
   X2 = Wire winding distance
   X3 = Wire guage

Since the maximum number of runs that could be
afforded timewise and costwise in this experiment
was n = 10, a 2^3 full factoral experiment
(involving n = 8 runs) could be afforded and so
was chosen to be run.  With an eye to the usual
monotonicity assumption in 2-level factorial
designs, the selected settings for the 3 factors
were as follows:

                                      -1         +1
                                     ----------------
   X1 = Number of wire turns :         90        180
   X2 = Wire winding distance:        .38        1.14
   X3 = Wire guage           :         40         48

The experiment was run (with the 8 settings
executed in random order). The following data
resulted:

<data>
   <add additional column: run sequence>
   <2 8 3 6 7 1 4 5>

5.6.1.2. Main Effects


-----Section 1: DEX Scatter Plot-----

<label the plot as a DEX Scatter Plot>
<impedance => sensitivity     everywhere>

...

<dex scatter plot>

Conclusions:

1. Important Factors:
   Factor 1 is clearly important ; when X1 = -1,
   all 4 sensitivities are low; when X1 = +1, all
   4 impedances are high.  Factor 2 is less
   important; the 4 sensitivities for X2 = -1 are
   slightly higher than the 4 sensitivities for X2
   = +1.  Factor 3 does not appear to be important
   at all--the sensitivity is about the same (on
   the average) regardless of the settings for X3.

2. Best Settings:
   The best settings depend on the definition of
   "best" for the response.  In this experiment,
   we are using the device as a detector, and so
   high sensitivities are desirable.  Given this,
   our first pass at best settings yields (X1 =
   +1, X2 = -1, X3 = either).

-----Section 2: DEX Mean Plot-----

<label the plot as a DEX Mean Plot>

...

<main effects plot>

Conclusions:

The main effects plot reaffirms the ordering of
the DEX scatter plot, but additional information is
gleaned because the eyeball distance between the
mean values gives an approximation to the least
squares estimate for the factor effects:

1. Important Factors:
   X1 (effect = large   : about 3 ohms)
   X2 (effect = moderate: about -1 ohm)
   X3 (effect = small   : about 1/4 ohm)

2. Best Settings:
   As before, choose those factor settings which
   (on the average) maximize the sensitivity:

      (X1,X2,X3) = (+,-,+)

-----Section 3: Interaction Effects Plot-----

In addition to the main effects, it is also
important to check for interaction
effects--especially 2-term interaction effects.
The DEX Mean Interaction Effects Plot is an
effective tool for this.  The 3 diagonal plots are
a replotting of the main effects-- they are plots
of the average response versus each of the 3
factors.  The off-diagonal plots have the average
response (vertically) versus the interaction (=
cross-product here) horizontally.  Thus for the
middle plot on the top, the horizontal axis is
X1*X2, where X1 takes on values -1 and +1, X2
takes on -1 and +1, and thus X1*x2 takes on values
-1 and +1.

The legend within each plot is the least squares
estimate of the factor (or interaction) effect; in
this orthogonal design case, the estimated effect
is identical to the difference of the mean values.

<matrix of plots>

Conclusions:

1. Important Factors:
   Noting the plots which have steepest lines (that is,
   largest effects), we note that:
      X1 is most important     : estimated effect = -3.1025;
      X2 is next most important: estimated effect = -.8675;
      X3 is relatively unimportant;
      All 3 2-term interactions are relatively unimporant.

2. Best Settings:
   As with the main effects plot, the best settings
   (to maximize the sensitivity) are

      (X1,X2,X3) = (+1,-1,+1)

    but with the X3 setting of +1 mattering little.

-----Section 4: Block Plots-----

Block plots are a useful adjunct to confirm the importance of
factors, to establish the robustness of main effect conclusions,
and to determine the existence of interactions.

The first plot below answers the question: Is factor 1 important
(robust over all 4 settings of X2 and X3)?
The second plot answers the question: Is factor 2 important
(robust over all 4 settings of X1 and X3)?
The third plot answers the question: Is factor 3 important
(robust over all 4 settings of X1 and X2)?

For block plots, it is the height of the bars which is important--
not the relative positioning of each bar; hence we focus on the
size and internals of the blocks--not "where" the blocks are one
relative to another.

<block plots>

Conclusions:

1. Relative Importance of Factors:
   All of the bar heights in plot 1 are bigger
   than the bar heights in plots 2 and 3--hence
   factor 1 is more important than factors 2
   and 3.

2. Statistical Significance
   In plot 1, looking at the levels within each
   bar, we note that the response for level 2
   is higher than level 1 in each and all of
   the 4 bars.  By chance this happens with
   probability 1/(2^4) = 1/16 = 6.25%.  Hence
   factor 1 is near-statistically significant
   at the 5% level.  Similarly, for plot 2,
   level 1 > level 2 for all 4 bars, hence
   factor 2 is near-statistically significant.
   For factor 3, there is not consistent
   ordering within all 4 bars and hence factor
   3 is not statistically significant.

3. Interactions:
   For factor 1, the 4 bars do not change
   height in any systematic way and hence there
   is no evidence of X1 interacting with either
   X2 or X3.  Similarly for factor 2.

-----Section 5: Estimating Effects-----

<text as you have it>

<yates output>

<omit "Interpretation" and have your text as is>
<I think "Interpetation" and "Conclusions" are close enough
to be confusing>

<in addition to the good stuff you  have explaining the Dataplot
output, I would suggest adding the following line...>

Of the 5 columns of output, the most important are the first
(which is the identifier), the second (which is the least squares
estimated effect = difference of means), and the last (which
is the residuals standard deviation for the cumualtive model--more on
this in the next section).

Conclusions:

1. Important Factors:
   X1 (Number of Turns)           : effect est. =  3.1025 ohms
   X2 (Winding Distance)          : effect est. = -0.8675 ohms
   X2*X3 (Number x Guage)         : effect est. =   .2975 ohms
   X1*X3 (Number x Guage)         : effect est. =   .2475 ohms
   X3 (Wire Guage)                : effect est. =   .2125 ohms
   X1*X2*X3 (Num. x Dist. x Guage): effect est. =   .1425 ohms
   X1*X2 (Number x Distance)      : effect est. =   .1275 ohms

-----Section 6: Modeling & Prediction Equations-----

In most cases of least squares fitting, the model
coefficient estimates for previously-added terms
change change depending on what was successively
added--hence the estimate for the X1 coefficient
might change depending on whether or not an X2
term was included in the model.  This is NOT the
case for when the design is orthogonal (as is this
2**3)--in such case, the estimates for the
previously-included terms are invariant of
successively added terms.  Hence the ranked list
of effect estimates in the above table
simultanously serves as the least squares
coefficient estimates for progressively more
complicated models.

The last column of the above table gives the
residual standard deviation for 8 possible
models--each one progressively more complicated.

At the top of the table, if none of the factors
are important, then the prediction equation
defaults to

      Yhat = the average = 2.65875

From the last column of the table, it is seen that
this simplest of all models has a residual
standard deviation (a measure of goodness of fit)
of 1.74106 ohms.  Finding a good-fitting model was
NOT one of the stated goals of this experiment,
but the determination of a good-fitting model is
"free" along with the rest of the analysis and so
it is included.

Conclusions:

From the last column of the above table, we
conclude the following: The prediction equation:

   Yhat = 2.65875 has a residual standard deviation
          has a residual standard deviation 1.74106 ohms
   Yhat = 2.65875 + 1/2 (3.1025*X1)
          has a residual standard deviation 0.57272 ohms
   Yhat = 2.65875 + 1/2 (3.1025*X1 - .8675*X2)
          has a residual standard deviation 0.30429 ohms
   Yhat = 2.65875 + 1/2 (3.1025*X1 - .8675*X2 + .2975*X2*X3)
          has a residual standard deviation 0.29750 ohms
   .
   .
   .
   and so forth.  The full-blown model:
   Yhat = 2.65875 + 1/2 (3.1025*X1 - .8675*X2 + .2975*X2*X3
          + ... + 0.12750*X1*X2) provides a perfect fit for the data.

-----Section 7: Intermediate Conclusions-----

Taking stock from all of the graphical and
quantitative analyses of previous sections, we
conclude that X1 (= number of turns) is the most
important engineering factor affecting
sensitivity, followed by X2 (= wire distance) as
next in importance, followed then by some
lesser-important interactions and X3 (= wire
guage).

Also, from the various analyses, we conclude that
the best design settings (on the average) for a
high-sensitivity detector is

   (X1,X2,X3) = (+,-,+)

that is

   (number of turns, wire distance, wire guage) = (180,.38,48)

Thus, in a very real sense, the analysis is
completed.  We have achieved the 2 most important
stated goals of the experiment:

   1. gaining insight as to what the most
      important factors are, and

   2. ascertaining as to what optimal production
      settings are.

On the other hand, more info can be squeezed from
the data and that is what the this and the
remaining sections address.

1. First of all, we focus on the problem of taking
   the ranked list of factors and objectively
   ascertaining which factors are "important"
   versus "unimportant".

2. In a parallel fashion with this, we use the
   above-derived subset of important factors to
   form a "final" prediction equation which is
   good (that is, having a sufficiently small
   residual standard deviation) while being
   parsimonious (having a small number of terms),
   as opposed to the already-derived default
   full-blown model which is perfect (having a
   residual standard deviation = 0, that is, the
   predicted values = the raw data) but unduly
   complicated (consisting of a constant + 7
   terms).

-----Section 8: Important Factors & Parsimonious Prediction Equation-----

The above 2 problems (important factors and
parsimonious model) will be handled in parallel.
Determination of one will yield the other.  In
regard to the "important factors", our immediate
goal is to take the full subset of 7 main effects
and interactions and extract a subset which we
will declare as "important", with the
complementary subset being "unimportant".  Six
criteria are utilized to define "important"
factors.  These 6 criteria are not all equally
important, nor will they yield identical subsets,
in which case a consensus subset or a weighted
consensus subset must be extracted.

Keep a factor and declare it to be in the
"important" subset if ...

   1. Effects                    : Engineering Significance
   2. Effects                    : Numerical Significance
   3. Effects                    : Statistical Significance
   4. Effects                    : Probability Plots
   5. Averages                   : Youden Plot
   6. Residual Standard Deviation: Engineering Significance
   7. Residual Standard Deviation: Statistical Significance

Thefirst 3 criterion are quantitative and us on
   effect estimates; the next criteria graphical
   and focus on effects estimates averages; the
   last 2 criteria are ntitative and focus on
   residual standard iations The last 2 criterion
   focus on the idual standard deviation

We discuss each of these 7 criteria in detail:

1. Effects: Engineering Significance
   <your formula>
   that is, declare a factor as "important" if
   the effect is > some a priori declared
   engineering difference.  This implies that the
   engineering has in fact stated as to what a
   minimum effect will be--oftentimes this is not
   the case.  In the absence of such, a good rough
   rule for the minimum engineering significant
   delta is to keep only those factors whose
   effect is greater then, say, 10% of the current
   production average.  In this case, let's say
   that the average detector has a sensitivity of
   2.5 ohms, then this would suggest that we would
   declare all factors whose effect is greater
   than 10% of 2.5 ohms = .25 ohm to be significant
   (from an engineering point of view).

   Based on the Minimum Engineering Significant
   Difference criterion, we thus conclude to keep
   2 terms: X1 (3.10250) and X2 (-.86750).

2. Effects: Numerical Significance
   <your formula>
   that is, declare a factor as "important" if
   one determines the largest effect and keep
   all factors in excess of 10% of that effect.
   We thus are keeping the "top 90%" of all
   effects.  This criterion is neither engineering
   nor statistical, but it does offer some
   additional numeric insight.  For the current
   case study, the largest effect is from factor 1
   (3.10250 ohms), and so 10% of that is .31 ohms,
   which suggests keeping all factors whose
   effects exceed .31 ohms.

   Based on the 90% Numerical criterion, we thus
   conclude to keep 3 terms:  X1 (3.10250) and X2 (-.86750)
   and X2*X3 (.29750).

3. Effects: Statistical Significance
   <your formula>
   that is, declare a factor as "important" if its
   effect is more than 2 standard deviations away
   from 0 (0, by definition, meaning "no effect").
   The difficulty with this is that in order to
   invoke this we need sigma = standard deviation
   of an observation.  This is problematic because

      1. the engineer may not know sigma;
      2. the experiment might not have
         replication, and so a model-free estimate
         of sigma is not obtainable;
      3. obtaining an estimate of sigma by
         assuming the sometimes- employed
         assumption of ignoring 3-term
         interactions and higher may be incorrect
         from an engineering point of view.

   For the current case study:
      1. the engineer did NOT know sigma;
      2. the design (a 2**3 full factorial) did NOT
         have replication;
      3. ignoring 3-term interactions and higher
         leads to an estimate of sigma based on
         omitting only a single term: the X1*X2*X3
         interaction.

   Thus for this current case study, if one
   assumes that the 3-term interaction is nil and
   hence represents a single drawing from a
   population centered at zero, then an estimate
   of the standard deviation of the effect is
   simply the effect itself (.1425).  2 such
   effect standard deviations is .2850--thus this
   rule becomes to keep all |effects| > .2850.  This
   results in the keeping of 3 terms: X1 (3.10250),
   X2 (-.86750), and X2*X3 (.29750).

4. Effects: Probability Plots
   1. Normal Probability Plot: Keep a factor as "important"
      if it is displaced off the line-through-zero
      on a normal probability plot of the
      effects.

   2. Halfnormal Probability Plot: Keep a factor as "important"
      if it is displaced off the
      line-near-zero on a halfnormal probability
      plot of the absolute value of effects.

   Both of these methods are based on the fact
   that the least-squares estimates of effects for
   these 2-level orthogonal designs are simply the
   difference of averages and so the central limit
   theorem suggests that (if no factor was
   important) the effects should have a normal
   distribution and the |effects| should have a
   halfnormal distribution.

   Since the halfnormal probability plot is only
   concerned with effect magnitudes as opposed to
   signed effects (which are subject to the
   vagaries of how the initial factor codings +1
   and -1 were assigned), then the halfnormal
   probability plot is to be preferred over the
   normal probability plot.

   For the case study at hand, both probability
   plots clearly show 2 factors displaced off the
   line, and from the third plot (with factor tags
   included), we see that those 2 factors are
   factor 1 and factor 2. All remaining 5 effects
   are behaving like random drawings from a normal
   distribution centered at zero, and so are
   deemed to be statistically insignificant.  In
   conclusion, this rule keeps 2 factors: X1
   (3.10250) and X2 (-.86750).

5. Effects: Youden Plot
   Keep a factor as "important" if it is displaced
   away from the central-tendancy "bunch" in a
   Youden plot of high and low averages.  By
   definition, a factor is important when its
   average response for the low (-1) setting is
   signficantly different from its average
   response for the high (+1) setting.
   Conversely, if the low and high averages are
   about the same, then what difference does it
   make which setting to use and so why would such
   a factor be considered important.  This fact in
   combination with the intrinsic benefits of the
   Youden plot for comparing pairs of items leads
   to the technique of generating a Youden plot of
   the low and high averages.

   <show Youden plot>

   For the case study at hand, the Youden plot
   clearly shows a cluster of points near the
   grand average (2.65875) with 2 displace points
   above (factor 1) and below (factor 2).  Based
   on the Youden plot, we thus conclude to keep 2
   factors: X1 (3.10250) and X2 (-.86750).

6. Residual Standard Deviation: Engineering Significance
   <sres > cutoff>
   that is, declare a factor as "important" until the
   cumulative model that includes the factor (and
   all larger factors) has a residual standard
   deviation smaller than an a prior
   engineer-specified cutoff.

   This criterion is different from the others in
   that it is model-focused.  In practice, this
   criterion states that starting with the largest
   effect, we cumulatively keep adding terms to
   the model and monitor how the residual standard
   deviation for each
   progressively-more-complicated model gets
   smaller and smaller.  At some point, the
   cumulative model will get complicated enough
   and comprehensive enough that the resulting
   residual standard deviation will drop below the
   pre-specified engineering cutoff for the
   residual standard deviation.  At such point, we
   stop adding terms and declare all of the
   model-included terms to be "important" and
   everything not in the model as "unimportant".

   This approach implies that the engineering has
   in fact thought through as to what a minimum
   residual standard deviation will be--in effect
   as to what the engineer can tolerate for the
   magnitude of the typical residual (= difference
   between the raw data and the predicted value
   from the model).  In other words, how good does
   the engineer want the prediction equation to
   be.  Unfortunately, this engineering
   specification has not always been formulated
   and so this criterion becomes moot.

   In the absence of such, a good rough rule for
   the minimum engineering residual standard
   deviation is to keep adding factors until the
   residual standard deviation just dips below,
   say, 5% of the current production average.  For
   this case study, let's say that the average
   detector has a sensitivity of 2.5 ohms, then
   this would suggest that we would keep adding
   factors to the model until the residual
   standard deviation first fell below 5% of 2.5
   ohms = .125 ohms.

   Based on the Minimum Residual Standard
   Deviation criterion, and by scanning the far
   right column of table ..., thus conclude from
   this criterion to keep 5 terms:

      X1       (with a cum. res. sd = .57272)
      X2       (with a cum. res. sd = .30429)
      X1*X2    (with a cum. res. sd = .26737)
      X1*X3    (with a cum. res. sd = .23341)
      X3       (with a cum. res. sd = .19121)
      X1*X2*X3 (with a cum. res. sd = .18031)
      X1*X2    (with a cum. res. sd = .00000)

   and thus we must include all terms in order to
   drive the residual standard deviation below
   .125.  Again, the 5% is a rough-and-ready rule
   which has no basis in engineering nor
   statistics, but rather "numerics".  Ideally,
   the engineer has a better cutoff for the
   residual standard deviation which is based on
   how well he/she wants the equation to peform in
   practice.  If such a number were available,
   then for this criterion and case study, we
   would end up with something less than the
   entire collection of factors.

7. Residual Standard Deviation: Statistical Significance
   <sres > cutoff>
   that is, declare a factor as "important" until
   the cumulative model that includes the factor
   (and all larger factors) has a residual
   standard deviation smaller than sigma--the
   standard deviation of an observation under
   replicated conditions.  In essence, we are
   allowing that we cannot demand of a model to
   fit any better than what we would get if we had
   replicated data; that is, we cannot demand that
   the residual standard deviation from any fitted
   model be any smaller than the (theoretical or
   actual) replication standard deviation.  We can
   drive the fitted standard deviation down (by
   adding terms) until it achieves a value close
   to sigma, but to drive it down further means
   that we are in effect fitting noise.

   In practice, this criterion may be difficult to apply because

      1. the engineer may not know sigma;
      2. the experiment might not have
         replication, and so a model-free estimate
         of sigma is not obtainable.

   For the current case study:
      1. the engineer did NOT know sigma;
      2. the design (a 2**3 full factorial) did
         NOT have replication.  The most common
         way of having replication in such designs
         would to have replicated center points,
         that is, replicates at the center of the
         cube, at (X1,X2,X3) = (0,0,0), that is,
         at (number of turns,winding distance,wire
         guage) = (135,.76,44).

   For the case study at hand,

      1. the experimentalist did not know sigma a
         priori; and

      2. the experiment did not have replication;

   hence estimates for sigma were not available.
   Thus for this current case, this criterion
   could NOT be used to yield a subset of
   "important" factors.


Conclusion for Important Factors:

In summary, the 7 criterion for specifying "important" factors
yielded the following:

   1. Effects: Eng. Signif.       : X1 X2
   2. Effects: Num. Signif.       : X1 X2 X2*X3
   3. Effects: Stat. Signif.      : X1 X2 X2*X3
   4. Effects: Prob. Plots        : X1 X2
   5. Averages: Youden Plot       : X1 X2
   6. Residual SD: Eng. Signif.   : all 7 terms
   7. Residual SD: Stat. Signif.  : not apply

Such conflicting results are common.  Arguably,
the 3 most important criteria (most important
listed first) are

   4. Effects: Prob. Plots        : X1 X2
   1. Effects: Eng. Signif.       : X1 X2
   6. Residual SD: Eng. Signif.   : all 7 terms

Scanning all of the above, we thus declare the
following consensus:

   1. Important Factors:
         X1 and X2
   2. Parsimonious Prediction Equation:
         Yhat = 2.65875 + 1/2 (3.10250*X1 - .86750*X2)
         (with a residual standard deviation of .30429 ohms)

-----Section 9: Validate the Fitted Model-----

The previous sections have let us to the model

   <equation>

The residual standard deviation of this model is
0.30429 ohms.

The next step is to validate the model.  The
primary method of model validation is graphical
residual analysis, that is, through an assortment
of plots of the differences between the observed
data y and the predicted value yhat from the
model.  Thus, for example, the design point
(-1,-1,-1) has an observed data point (from table
...) of y = 1.70, while the predicted value from
the above fitted model for this design point is

   yhat = 2.65875 + 0.5(3.10250*(-1) - 0.86750*(-1))
   yhat = ...

which leads to the residual

   res = y - yhat = ...

If the model fits well, then yhat should be near y
for all 8 design points.  Hence the 8 residuals
shoukld all be near zero.  The 8 predicted values
and residuals for this data and model are:

   X1   X2   X3  Observed Predicted  Residual
----------------------------------------------
   -1   -1   -1    1.70    1.54125    0.15875
   +1   -1   -1    4.57    4.64375   -0.07375
   -1   +1   -1    0.55    0.67375   -0.12375
   +1   +1   -1    3.39    3.77625   -0.38625
   -1   -1   +1    1.51    1.54125   -0.03125
   +1   -1   +1    4.59    4.64375   -0.05375
   -1   +1   +1    0.67    0.67375   -0.00375
   +1   +1   +1    4.29    3.77625    0.51375

What is the magnitude of the typical residual?
There are several ways to compute this, but the
statistically optimal measure is the residual
standard deviation:

   sres = sqrt(sum of squared residuals)/(n-p)

where for this case the number of observations n =
8 and the number of fitted parameters p = 3. We
know from previous calculations and as noted at
the beginning of this section that sres = .30429.

If the prediction equation is adequate as is, then
the residuals from that equation should behave
like random drawings (typically from a normal
distribution), and such residuals should--since
presumably random--have no structural relationship
with any factor--including any and all potential
terms (X1, X2, X3, X1*X2, X1*X3, X2*X3, X1*X2*X3).

Further, if the model is adequate and complete,
then the residuals should have no structural
relationship with ANY other variables that may
have been recorded (for example, run
sequence--time-- which is really serving as a
surrogate for any physical or environmental
variable correlated with time).  Ideally, all such
residual scatter plots should appear
structureless.  Any scatter plot which exhibits
structure suggests that that factor should have
been formally included as part of the prediction
equation.

Validating the prediction equation thus means that
we do a final check as to whether any other
variables may have been inadvertantly left out of
the prediction equation, including variables
drifting with time.

The graphical residual analysis thus consists of
scatter plots of the residuals versus all 3 main
factors and 4 interactions (all such plots should
be structureless), a scatter plot of the residuals
versus run sequence (which also should be
structureless), and a normal probability plot of
the residuals (which should be near-linear).  We
present such plots below.

<7 plots + 2 plots>

Conclusions

1. Main Effects & Interactions:
   The X1 and X2 scatter plots are "flat" (as they
   must be since X1 and X2 were explicitly
   included in the model).  The X3 plot shows some
   structure as does the X1*X3, the X2*X3, and the
   X1*X2*X3 plots.  The X1*X2 shows little
   structure.  The net effect is that the relative
   ordering of these scatter plots is very much in
   agreement (again, as they must be) with the
   relative ordering of the "unimportant" factors
   given on lines 3 to 7 of the ...  table.  From
   the ...  and from the X2*X3 plot, it is seen
   that the next most influential term to be added
   to the model would be X2*X3.  In effect, these
   plots offer a higher-resolution confirmation of
   the ordering that was in table ...  On the
   other hand, none of these other factors
   "passed" the 7 criteria as given in the
   previous section and so these
   factors--suggestively influential as they might
   be-- are still not influential enough to be
   added to the model.

2. Time Drift:
   The run sequence scatter plot is random.  Hence
   there does not appear to be a drift, either
   from time, or any factor (e.g., temperature,
   humidity, pressure, etc.) possibly correlated
   with time.

3. Normality:
   The normal probability plot of the 8 residuals
   has some curvature which suggests that
   additional terms be added.  On the other hand,
   the correlation coeffient of the 8 ordered
   residuals and the 8 theoretical normal N(0,1)
   order statistic medians (which define the 2
   axes of the plot) has the value ..., which is
   well within acceptable (5%) limits of the
   normal probability plot correlation coefficient
   test for normality [ref].  In effect, therefore
   the plot is not so non-linear as to reject
   normality.

In toto, therefore, we end up accepting the model

<equation>

as a parsimonious, but good representation of the
sensitivity phenomenon under study.

-----Section 10: Using the Fitted Model-----

Although deriving such a model was not the primary
purpose of the study, it does have 2 benefits in
terms of additional insight:

   1. Global prediction
   2. Global determination of best settings.

1. Global prediction:

How does one predict the response at points other
than those used in the experiment?  The prediction
equation yields good results at all 8 combinations
of the coded -1 and +1 values for the 3 factors:

   Number of turns = 90 and 180
   Winding distance = .38 and 1.14
   Wire guage = 40 and 48

What, however, would one the expect the detector to yield
at uncoded target settings of, say,

   Number of turns = 150
   Winding distance = .50
   Wire guage = 46

The prediction equation allows us to answer this question.
Based on the fitted equation, we first translate
the target settings into coded target settings (on the usual
-1 to +1 scale) as follows:

   coded target setting = -1 + 2*p

where p = (uncoded target setting-low)/(high-low).  Hence the
coded target values are

   X1 = -1 + 2*(150-90)/(180-90) = .333333
   X2 = -1 + 2*(.50-.38)/(1.14-.38) = -.684211
   X3 = -1 + 2*(47-40)/(48-40) = .75000

Thus the uncoded target setting of

   (Number of turns,Winding distance,Wire guage) = (150,.50,47)

translates into the coded target setting of

   (X1,X2,X3) = (.333333,-.684211,.750000).

Inserting these coded values into the fitted
equation yields the desired predicted value of

   yhat = 2.65875 + 0.5(3.10250*(.333333) - 0.86750*(-.684211))
        = 3.47261

The above procedure can be applied for any
values of turns, distance, and guage-- subject to
the usual cautions that equations that are good
near data point vertices may not necessarily be
good everywhere in 3-space.  Interpolation is
typically a bit safer than extrapolation, but there
is never any assurance in either case.

With respect to interpolation,
one would feel more comfortable about
prediction (as in our example) if

   1. additional data had been collected at the
      center points and

   2. the center point data turned out to be in
      good agreement with predicted values at the
      center point based on the fitted model.

In our case, we had no such data and so the
simple truth is that the use of the prediction
equation even for interpolation is assuming something in which the data
set as given is not capable of proving one way or
the other.  Given that reality, the purpose of the above was to
demonstrate how one may cautiously generate
predicted values which go well beyond our limited
original data set of 8 points.

2. Global Determination of Best Settings

Given the existence of a prediction equation, what
(additional) information does that convey to us in
terms of "best settings" when such best settings
may lie OUTSIDE of the domain in which the
experiment was carried out?  In particular,
suppose the truth of the matter is that the
nature's optimal settings of

   (number of turns, distance, guage)

is in fact outside the test region:

   X1 = Number of wire turns :         90        180
   X2 = Wire winding distance:        .38        1.14
   X3 = Wire guage           :         40         48

Can we detect such an extra-domain optimal point?
The answer is this is a cautious yes and clues to
such are most easily arrived at via contour plots.
In particular, the contour plot of choice involves
the 2 dominant factors in the analysis.

IN our case study (see table xxx) the 2 such factors
were X1 and X2.
The X1 (horizontal) and X2 (vertical) contour plot is given below.

<show contour plot>

Construction-wise,
the 4 vertex values of this imbedded 2**2 full
factorial design involving X1 and X2 were obtained
from the raw data (see xx) by simple averaging.
Thus the value on the contour plot at X1 = -1 and
X2 = -1 was gotten by referencing the raw data
table and noting all (X1 = -1, X2 = -1) data;
there were 2 such data values: Y = 1.51  and Y = 1.70.
Averaging these 2 values together yields
the 1.6 which in turn appears on the plot.  We proceed
similarly for the other 3 vertices.  To draw the
contour lines, we make use of the local 2-term equation
in X1 and X2 only, namely:

   Yhat = mu + 0.5(B1*X1 + B2*X2 + B12*X1*X2)
   Yhat = 2.65875 + 0.5

where the coefficients B. were drawn from the
least squares estimates as given in table ...  .
From the raw data table (see xx), we note that the
8 data values ranged from ...  to ..., and so
contour lines were drawn from ...  to ..., at
increments of ...  .

From the contour plot, we may derive the following
3 pieces of information:

   1. Interaction significance;
   2. Best (data) setting for these 2 dominant factors;
   3. Best (non-data) settings.

1. Interaction Significance:
   Note the appearance of the contour plot--if the
   contour curves are linear, then that implies
   that the interaction term is insignificant; if
   the contour curves have considerable curvature,
   than that implies that the interaction term is
   large and important.  In our case, the contour
   curves are linear and so we conclude that the
   X1*X2 term is insignificant-- that is, we
   conclude that the X1 effect (= ...  units) is
   valid irrespective of the settings of X2.

2. Best (data) Settings:
   To determine the best (data) setting (that is,
   the best factor settings for the already-run
   experiment), then we first must define what
   "best" means.  As before, for this case study,
   "best" here means to MAXIMIZE (rather than
   minimize or hit a target) the response.  Hence
   from the contour plot, we determine the best
   settings for the 2 dominant factors by simply
   scanning the 4 vertices and choosing the vertex
   with the LARGEST value (= average response).
   In this case, it is (X1 = +1, X2 = +1).  Hence
   the best settings for the 2 dominant factors
   based on the experiment that was already run is

      (X1 = +1, X2 = +1).

   As for factor X3, the contour plot provides no
   best setting info, and so we would resort to
   other tools: the main effects plot, the
   interaction effects matrix, or the ordered data
   (see next section) to determine optimal X3
   settings.

3. Best (non-data) Settings:
   To determine the best (non-data) setting (that
   is, the best setting even beyond the settings
   utilized in the already-run experiment), we do
   the following 6 steps:

      1. Note Best Current Value;
      2. Specify Target Future Value;
      3. Note Target Contour Curve;
      4. Determine Closest Point on Target Contour Curve;
      5. Translate Optimal Settings to Original Factor Units;
      6. Do Confirmation Run or Confirmation Experiment.

   We now discuss each of the above 6 steps in detail:

      1. Note Best Current Value:
         We first go to the best data setting;
         that is, we use as our starting point
         that vertex (out of the 4) with the
         highest average response.  For the case
         study at hand, that optimal point is
         (X1 = +1, X2 = +1).

      2. Specify Target Future Value:
         We mentally specify as to what value of
         the response (above and beyond our best
         data vertex) that we would like to see if we
         had the luxury of sampling elsewhere
         in the factor domains.
         For our case study, we can already
         achieve (on the average) ..., and so let us suppose we
         choose to increase our response to ... --
         we thus specify ... as the "target future value".

      3. Note Target Contour Curve:
         We note the plot contour curve which comes
         closest to the target future value.  Our
         next best setting is thus going to be
         somewhere along the Y = ...  contour
         curve.
         In our case, with the target future value at ...,
         and with the contour curves at ..., ..., ..., to ...
         we thus choose contour curve Y = ... as the target curve.

      4. Determine Closest Point on Target Contour Curve
         Since the contour curves were based on the
         data and since the validity of the
         contour curves becomes more and more
         suspect as we move away from the data
         vertices, then the strategy is to choose
         that value on the target contour curve which is
         closest to the optimal data vertex.  In
         practice this is achieved by starting at
         this best vertex and proceeding in the
         direction of the path of steepest ascent
         (noting the nature and curvature of the
         local contour curves) until we
         (perpendicularly) intersect the target
         contour curve.  The interesection point
         becomes our best "estimate" for the X1 and X2
         settings that will yield the desired
         future value.  This is the point where we
         would like to conduct our next
         experiment, or at least run confirmation
         runs to determine if this best-estimate
         point in fact yields the quality of
         response that we desire.

      5. Translate Optimal Settings to Original Factor Units
         At this intersection point, note the
         coded settings for factors X1 and X2.
         Since the original 4 vertices were coded
         values of +-1 in each of the 2 factors
         under study, it is quite common for the
         best new settings to be somewhere in the
         region +- 1.25 to +- 2. Determine exactly
         such best settings in coded units, and
         then linearly translate back to the best
         settings in the original factor units.
         For example, in our case, if the optimal
         new setting for X1 (= number of wire
         turns) was X1 = +2, and since

            X1 = -1 corresponds to  90 turns, and
            X1 = +1 corresponds to 180 turns, then
            X1 = +2 corresponds to 225 turns.

         Similarly, if X2 (winding distance) = +2
         is optimal, then since

            X2 = -1 corresponds to  .38, and
            X2 = +1 corresponds to 1.14, this implies
            X2 = +2 corresponds to 1.52.

      6. Do Confirmation Run or Confirmation Experiment
         Test the above best setting by
         additionally collecting data at this new
         point, or running a 2-level factorial
         experiment (centered at this point) in
         the vicinity of this new point (and
         repeat the entire 2-level analysis
         again).

-----Section 11: Best Settings (Summary)-----

At this point it is good to review the conclusions
regarding goal 2 (determine best settings) for
this case study.  If we are trying to determine
settings of the factors which maximize the
response, as in the this eddy current case study,
then there are several approaches which lay claim
title to yielding "best" settings:

   1. Average-based
   2. Model-based
   3. Data-based

Methods 1 and 3 yield answers within the current
-1 and +1 domain of the data.  Method 2 yields an
answer that frequently extrapolates beyond the
domain of the executed experiment design.

This eddy current case study was drawn from an R & D
environment, and so heavier emphasis was placed on
understanding and insight (that is, what factors
are important) than with best settings.  On the
other hand, in a production environment, the both
goals are still in play, but the emphasis goes
slightly heavier to the best settings goal as
opposed to the important factors goal.  In a
general engineering environment (either R & D or
production), the above 2 goals frequently have, in
practice, much greater interest to the engineer
that the third enumerated goal (good
model/prediction equation).

What then, are the final recommended best settings
for this case study, and how do they compare with
one another?

1. Average-based:
   From the main effects plot (or equivalently
   from the signs of the estimated effects), it
   was seen that the comparison of the average
   response for each individual factor at the -1
   and + levels, and the choosing of the setting
   that yielded the larger average response
   yielded

      (x1 = +1, X2 = +1, x3 = +1)

   From step 1 of the contour plot analysis, we
   saw that the best of the 4 vertices for the 2
   dominant factors (X2 and X3):

      (X2 = -1, X3 + 1)

   which agrees in part with the above seetings
   but leaves the question of X1 undone.

2. Model-based:
   Also from the contour plot, we were able to
   draw 2-factor model-based contour curves which
   yielded

      (X1 = ?, X2 = -2, X3 = +2)

   as the global estimate of the best setting
   (with the setting of X1 left undetermined).

3. Data-based
   Finally, a different, assumption-free, and
   simple (= no estimation and no modeling)
   data-based estimate of the best settings is
   always of interest to pursue.  If one takes the
   decidedly anti-statistical attitude that all of
   the previously-described statistical estimation
   and statistical modeling is nice but is a bit
   artificial and indirect, and that the only
   thing real is the 8 data values on hand, then
   an argument could be forwarded that a simple
   sort of the 8 data values (carrying along the 3
   factor settings) should yield interesting
   information and optimal values all by itself.
   Proceeding in this fashion, such a simple sort
   (smallest to largest) of the data was carried
   out, yielding:

   <show sorted data:>
       Y         X1          X2          X3
     Probe     Number      Winding      Wire
   Impedance  of Turns     Distance     Guage
   -------------------------------------------
      0.55       -1          +1          -1
      0.67       -1          +1          +1
      1.51       -1          -1          +1
      1.70       -1          -1          -1
      3.39       +1          +1          -1
      4.29       +1          +1          +1
      4.57       +1          -1          -1
      4.59       +1          -1          +1

   The following conclusions can be drawn from the sort:

      1. Best Settings:
         The largest response is 4.59 and so the
         corresponding factor settings:

            (X1 = +1, X2 = -1, X3 = +1)

         can validly lay valid claim to being the
         best settings.  Note interestingly enough
         that worst data point (0.55) has the
         exact opposite signs:

            (X1 = -1, X2 = +1, X3 = -1)

         which further reinforces our choice for
         the (+1,-1,+1) best setting.  That
         happened in this case; it does not
         happen in most cases.

      2. Important factors
         Note that the largest value (4.59) has
         settings (+,-,+).  Note that the next
         largest (and nearly identical) value is
         (4.57) and it has settings (+,-,-).  The
         2 common settings for the 2 largest
         values are X1 (+) and X2 (-) which
         supports the fact that they are indeed
         important factors, while X3 is relegated
         to being an unimportant factor-- as per
         agreement with the extensive previous
         statistical analyses.

         Note also that in the sorted table, the 4
         largest values all had X1 = +1 and the 4
         smallest values had X1 = -1, thus
         reaffirming that factor X1 is in fact the
         most important factor.  Further, within
         the largest 4 responses with X1 = +1,
         note that X2 is +1 for 2 largest values
         and -1 for the next 2 values.  Similarly,
         for the 4 smallest values, the same
         pattern for X2 (2 +'s followed by 2 -'s)
         persist, thus supporting the prior
         conclusion that X2 is the second most
         important factor.

   4. Comparing Best Settings: Average-based vs. Data-based
      Based on the averages, the derived best
      setting was (+ -,+).  Based on the sorted
      data, the derived best setting was (+,-,+).
      Fortuitously in this case, the 2 different
      approaches yielded the same conclusion, and
      so our final unambiguous estimate for the
      best setting within the domain of the data is

         (X1 = +1, X2 = -1, X3 = +1)

      Because 2 different methods converged to the
      same conclusion, our confidence in the
      validity of that conclusion is greatly
      enhanced.

      In other case studies, the 2 approaches may
      not agree in all of the settings.  In this
      author's experience, however, the 2 methods
      will ALWAYS agree for the important factors
      (factors X3 and X2 here) and disagreement
      (if any) shows up in the unimportant
      factors.  In such cases, the recommendation
      course of action is to

         1. Important factors:
            Use the common best settings as given
            by either and both approaches.

         2. Unimportant Factors:
            1. If there is a difference in cost or
               set-up time for the unimportant
               factors, then use the least
               expensive or most hassle-free
               settings for these factors.

            2. If there is no difference in factor
               cost or set-up time, and the
               analyst must decide on the
               average-based settings for the
               unimportant factors, or the
               data-based settings for the
               unimportant factors, then this
               author would recommend using the
               average-based for the following 2
               reasons:

                  1.The scope is broader for
                    average-based.  The settings
                    depend on averages across 2 or
                    more observations, as opposed
                    to having the noisiness of
                    just a single observation in
                    the data-based method.

                  2.The settings for the
                    unimportant factors will
                    always exist for full
                    factorial designs (as is our
                    present case study) but for
                    other studies using a
                    fractional factorial design,
                    the use of average-based
                    estimates for the settings of
                    unimportant factors can yield
                    settings for which data may
                    not yet exist, and so is a
                    step beyond the settings based
                    just on the sorted data.

